\section*{문제 (본고사 97학년도)}

한 변의 길이가 $3r$ cm인 정육면체 모양의 그릇에 물이 반 정도 차 있다.
이 그릇에 반지름의 길이가 $r$ cm인 공이 수면에 수직인 방향으로
$1$ cm/초의 속력으로 잠긴다고 하자.
공의 반이 물에 잠기는 순간의 수면의 상승속력을 구하라.

\section*{풀이 A — 수면 기준 속도 $\left(\dfrac{d\xi}{dt}=1\right)$}

\subsection*{부피 분해}

$V = C + W$

\noindent $C$: 최초 물의 부피 (상수) \quad $W$: 물에 잠긴 구의 부피

\subsection*{잠긴 구의 부피 $W$}

$\xi$ : 수면 아래로 잠긴 깊이. 구의 잠긴 부분을 회전체로 보면:
\[
  W = \pi \int_{r-\xi}^{r} (r^2 - y^2)\,dy
\]

\subsection*{전체 부피 방정식}

정육면체 밑넓이 $= (3r)^2 = 9r^2$ 이므로:
\[
  9r^2\, h = C + \pi \int_{r-\xi}^{r} (r^2 - u^2)\,du
\]

\subsection*{$t$에 대해 미분}

\[
  9r^2 \frac{dh}{dt} = \pi\,\frac{d\xi}{dt}(2r\xi - \xi^2)
\]

\subsection*{$\dfrac{d\xi}{dt} = 1$ 대입}

\[
  \frac{dh}{dt} = \frac{\pi(2r\xi - \xi^2)}{9r^2}
\]

\subsection*{공의 반이 잠기는 순간 $(\xi = r)$}

\[
  \frac{dh}{dt}
    = \frac{\pi(2r^2 - r^2)}{9r^2}
    = \frac{\pi r^2}{9r^2}
\]

\[
  \boxed{\dfrac{dh}{dt} = \dfrac{\pi}{9} \text{ cm/초}}
\]

\section*{풀이 B — 밑바닥 기준 속도 $\left(\dfrac{d\ell}{dt}=-1\right)$}

\subsection*{변수 정의}

$\ell$ : 그릇 밑바닥에서 구의 아랫부분까지의 거리.
구가 1 cm/초로 가라앉으므로 $\dfrac{d\ell}{dt} = -1$.

$\xi$ 와 $\ell$ 의 관계:
\[
  \xi = h - \ell, \qquad
  \frac{d\xi}{dt} = \frac{dh}{dt} - \frac{d\ell}{dt} = \frac{dh}{dt} + 1
\]

\subsection*{미분 방정식에 대입}

풀이 A의 미분 결과에 $\dfrac{d\xi}{dt} = \dfrac{dh}{dt}+1$ 을 넣으면:
\[
  9r^2 \frac{dh}{dt} = \pi(2r\xi - \xi^2)\!\left(\frac{dh}{dt} + 1\right)
\]

\subsection*{공의 반이 잠기는 순간 $(\xi = r)$}

\begin{align*}
  9r^2 \frac{dh}{dt} &= \pi r^2\!\left(\frac{dh}{dt} + 1\right) \\
  9r^2 \frac{dh}{dt} &= \pi r^2 \frac{dh}{dt} + \pi r^2 \\
  (9 - \pi)\,r^2 \frac{dh}{dt} &= \pi r^2
\end{align*}

\[
  \boxed{\dfrac{dh}{dt} = \dfrac{\pi}{9 - \pi} \text{ cm/초}}
\]
